% ============================================================
% resumen.tex — Resumen/Abstract del paper
% ============================================================
\begin{abstract}
El fraude en transacciones con tarjetas de crédito representa una amenaza
creciente para el sistema financiero global. Este trabajo presenta una
solución basada en aprendizaje automático para la detección automática de
transacciones fraudulentas, empleando la metodología CRISP-DM sobre el
dataset Credit Card Fraud Detection disponible en Kaggle, que contiene
284,807 transacciones de titulares europeos con un desbalance severo de
clases (0.17\% de fraudes). Se evaluaron cuatro algoritmos de
clasificación: Regresión Logística, Random Forest, XGBoost y Gradient
Boosting. El manejo del desbalance se realizó mediante SMOTE. Los
resultados demuestran que [completar con mejor modelo] alcanzó el mayor
ROC-AUC de [valor], con un Recall de [valor] para la clase fraude,
evidenciando la viabilidad de los modelos propuestos para su integración
en sistemas de detección en tiempo real.

\textbf{Palabras clave:} fraude, clasificación, SMOTE, Random Forest,
XGBoost, desbalance de clases, CRISP-DM.
\end{abstract}
